% =====================================================================================================
% zitex in a research paper: a demo with annotations
%
%   Build it (in this folder; `zitex setup` and `zitex extract` already made everything it needs):
%       xelatex main.tex          (run twice for the cross-references)
%
%   How the pieces fit together
%       chars.yaml + lexicon.yaml  --zitex extract-->  zitex-out/<char>-<style>[-<lang>][@reading].tex
%       main.tex  --\zi... commands-->  \input of exactly those files   (zitex.sty is a thin wrapper)
%   So LaTeX never parses a decomposition: Python wrote the TikZ / amsmath snippet beforehand.
%   A character with no snippet shows as a red [字?] and gives a "Missing zitex-out/..." warning in
%   main.log. Fix: add it to chars.yaml (or run `zitex extract -c ...` again) and re-run xelatex.
% =====================================================================================================
\documentclass[11pt]{article}
\usepackage[margin=2.2cm]{geometry}
\usepackage{zitex}                       % ANNOTATION: loads xeCJK, fontspec, amsmath, tikz, graphicx; defines \zi \zimath \zimatrix \zitree \ziformula
\setCJKmainfont{Noto Serif CJK SC}       % the primary font of fonts.yaml; \setCJKmainfont here overrides the package default
\usepackage{booktabs}
\usepackage{caption}
\usepackage{hyperref}

\captionsetup{font=small,labelfont=bf}

% ----- helpers used only to annotate this demo (not part of zitex) ---------------------------------
\newcommand{\cmd}[1]{\texttt{\textbackslash #1}}               % a command name in running text
\newenvironment{syntax}{\par\smallskip\begin{minipage}{\linewidth}\footnotesize\setlength{\parindent}{0pt}\noindent\ignorespaces}
  {\end{minipage}\par\smallskip}
\newcommand{\shape}[1]{{\small #1}}
\newcommand{\note}[1]{\par\noindent\hangindent=1.2em\hangafter=1{\small\textbf{Syntax.}\ #1}\par\smallskip}

\title{Chinese Characters as Reactions:\\ Typesetting Decompositions with \texttt{zitex}}
\author{Demo author\thanks{Every figure, equation and formula below comes from \cmd{zi}-commands. The text around them is ordinary \LaTeX.}}
\date{}

\begin{document}
\maketitle

\begin{abstract}
We write the decomposition of a Chinese character in a few lines of data, and \texttt{zitex} turns it into
a tree, a nested Zi-Matrix, a one-line formula and a typeset equation. This file is both a short paper and a
guide: each block of markup is followed by a note on its syntax, and the features it shows.
\end{abstract}

% =====================================================================================================
\section{Introduction}
% =====================================================================================================
A character can be read as the product of a reaction, for example $\text{亻} + \text{弗} \rightarrow \text{佛}$.
Writing that by hand in every paper is slow, and the layout of the parts (left--right, top--bottom, enclosing) is
lost. With \texttt{zitex} the same thing in the middle of a sentence is \ziformula{佛}\unskip, and the layout is kept in
the matrix of Figure~\ref{fig:matrix}.

\note{\cmd{ziformula}\texttt{\{字\}} prints \emph{parts $\rightarrow$ character (pinyin, gloss)} on one line, so it works
inside a sentence, a caption or a table cell. It is the only style that is not a block.}

\textbf{The data.} A character is a short YAML entry: its shape in prefix form (for 佛, 亻 at \texttt{L} and 弗 at \texttt{M}), plus the position of each
part in its level. The positions are the eleven Zi-Matrix codes
\texttt{L R M U D Mi Mo LU LD RU RD}, where \texttt{M} is the centre of gravity and a Kangxi radical is peripheral.
This paper uses three characters, taken from a database with \texttt{zitex extract}:

\begin{center}
\small
\begin{tabular}{@{}cllll@{}}
\toprule
Character & Pinyin & Meaning & First-level pattern & First-level parts \\
\midrule
器 & qì  & device         & \texttt{LU-RU-M-LD-RD} & \shape{口 口 犬 口 口} \\
藻 & zǎo & aquatic grasses & \texttt{U-D}           & \shape{艹 + 澡} \\
佛 & fó  & Buddha         & \texttt{L-M}           & \shape{亻 + 弗} \\
\bottomrule
\end{tabular}
\end{center}

\note{The long form \cmd{zi}\texttt{[style=formula]\{器\}} and the short form \cmd{ziformula}\texttt{\{器\}} are the same.
Every style has a shortcut: \cmd{zitree}, \cmd{zimatrix}, \cmd{ziformula}, \cmd{zimath}. The parts
are read from \texttt{chars.yaml}, and the table is typed by hand.}

% =====================================================================================================
\section{Four views of one decomposition}
% =====================================================================================================
\subsection{Equation: \texttt{\textbackslash zimath}}
The equation view is the most compact. Pinyin sits above each part, and the gloss below.

\zimath{藻}

\note{\cmd{zimath}\texttt{\{字\}} is a display equation made only of \cmd{text}, \cmd{overset}, \cmd{underbrace}
and \cmd{xrightarrow}, so the same snippet works in MathJax. A part that has parts of its own is written as a
bracketed group, which is how 藻 shows 澡 inside it (see the tree and the matrix below). The \texttt{lang} key picks the gloss
language, as in \cmd{zimath}\texttt{[lang=fr]\{悔\}}, and needs that language in the lexicon.}

\subsection{Matrix: \texttt{\textbackslash zimatrix}}
The matrix keeps the \emph{geometry}: parts are boxes placed where the position codes say.

\begin{figure}[ht]
  \centering
  \zimatrix{器}\qquad\zimatrix{佛}
  \caption{Zi-Matrices of 器 (left: four 口 at the corners and 犬 in the centre) and 佛 (right).}
  \label{fig:matrix}
\end{figure}

\note{\cmd{zimatrix}\texttt{\{字\}} draws with TikZ and is an ordinary box, so it goes in a \texttt{figure}, a table cell,
or a \cmd{resizebox}. Two matrices fit side by side here because each is narrow. The default
text size is set by the snippet. To change it, scale the box: \cmd{resizebox}\texttt{\{0.8\textbackslash linewidth\}\{!\}\{\textbackslash zimatrix\{器\}\}}.}

\subsection{Tree: \texttt{\textbackslash zitree}}
The tree shows how deep a decomposition goes. 藻 has four levels, because 澡 and 喿 are themselves decomposed.

\begin{figure}[ht]
  \centering
  \resizebox{0.9\linewidth}{!}{\zitree{藻}}
  \caption{The decomposition tree of 藻. Leaves are the parts that are not decomposed further.}
  \label{fig:tree}
\end{figure}

\note{A tree can be wide, so Figure~\ref{fig:tree} is wrapped in \cmd{resizebox}. The wrapper only scales,
the snippet itself is unchanged. \cmd{zitree} accepts the same keys as the others:
\texttt{style}, \texttt{lang}, \texttt{reading}.}

\subsection{Choosing the style as an option}
All four commands are shortcuts for one: \cmd{zi}\texttt{[style=...]}.

\begin{center}
\small
\begin{tabular}{@{}lll@{}}
\toprule
Command & Output for 佛 & Use it for \\
\midrule
\cmd{zimath}\texttt{\{佛\}}    & an equation (display)        & the main text, one idea per equation \\
\cmd{zimatrix}\texttt{\{佛\}}  & nested boxes                 & figures, showing the layout \\
\cmd{zitree}\texttt{\{佛\}}    & a tree                       & figures, showing the depth \\
\cmd{ziformula}\texttt{\{佛\}} & \ziformula{佛}\unskip              & a sentence, a caption, a table \\
\bottomrule
\end{tabular}
\end{center}

% =====================================================================================================
\section{Results: a small gallery}
% =====================================================================================================
The same three characters as a matrix and as a formula:

\begin{table}[ht]
  \centering
  \caption{One row per character. Cells hold \cmd{zi} output directly.}
  \label{tab:gallery}
  \small
  \begin{tabular}{@{}c c c@{}}
    \toprule
    & Matrix & Formula \\
    \midrule
    器 & \resizebox{0.40\linewidth}{!}{\zimatrix{器}} & \ziformula{器} \\[4pt]
    藻 & \resizebox{0.40\linewidth}{!}{\zimatrix{藻}} & \ziformula{藻} \\[4pt]
    佛 & \resizebox{0.40\linewidth}{!}{\zimatrix{佛}} & \ziformula{佛} \\
    \bottomrule
  \end{tabular}
\end{table}

\note{Block styles work in a \texttt{tabular} as long as the column is wide enough, which is why the matrices are
scaled with \cmd{resizebox}. \cmd{ref}, \cmd{label} and \cmd{caption} behave as for any figure or table, so a paper can
refer to Table~\ref{tab:gallery}, Figure~\ref{fig:matrix} and Figure~\ref{fig:tree} in the usual way.}

\textbf{Observation.} In 器, a dog (犬) sits inside four mouths (口). In 藻 the grass radical (艹) is on top and the water radical
(氵) is a part of 澡, which is why the equation reads \ziformula{藻}\unskip{} and the tree goes deeper (Figure~\ref{fig:tree}).
This is a description of the data, not a claim about etymology.

% =====================================================================================================
\section{Features you will want in a real paper}
% =====================================================================================================
\subsection{Language and reading}
Equations and formulas print glosses, so they depend on a language. Set a default for the document or override it in one call:

\begin{syntax}
\begin{verbatim}
\zitexsetlang{fr}                % glosses in French from here on
\zimath[lang=en]{藻}             % one call in English
\zitexsetreading{dictionary}     % use a named reading
\zimath[reading=author]{虹}      % one call with another reading
\end{verbatim}
\end{syntax}

\note{A \emph{reading} is a named interpretation of the roles of the parts (semantic, phonetic, force), so two schools of
thought can sit in the same data and a paper can choose. In the file system they appear as
\texttt{<char>-<style>[-<lang>][@reading].tex}. These lines are not run here, because the lexicon of this folder has
English only and no readings.}

\subsection{Where the snippets are}
\begin{syntax}
\begin{verbatim}
\zitexsetpath{zitex-out}         % the default: a folder next to main.tex
\end{verbatim}
\end{syntax}

\note{To share one folder between several papers, point every paper to it, and run \texttt{zitex extract -o} on that folder.}

\subsection{Fonts}
The \texttt{zitex} package has a font fallback: when the primary font has no glyph for a part (rare characters, ancient components), a fallback font,
an SVG image or an alias is used, as set in \texttt{fonts.yaml}. A part that nothing can draw prints a boxed code point such as
\texttt{U+XXXX} and is listed by \texttt{zitex fonts}.

% =====================================================================================================
\section{Reproducing this paper}
% =====================================================================================================
\begin{syntax}
\begin{verbatim}
zitex setup   -o work/                         # zitex.sty, fonts.yaml, glyphs/
zitex extract -c 器,藻,佛 -o work/            # chars.yaml, lexicon.yaml, main.tex, zitex-out/
cd work && xelatex main.tex                    # the PDF
\end{verbatim}
\end{syntax}

\note{\texttt{extract} writes a minimal \texttt{main.tex} and never replaces one that exists unless you pass \texttt{--force}.
To use this file instead, copy it over \texttt{main.tex}, and add any new character to the \texttt{-c} list. If a
character is missing, the red mark and the log line \texttt{Missing zitex-out/...} say which.}

\medskip
\noindent\textbf{Summary of the syntax.}
\begin{center}
\small
\begin{tabular}{@{}p{0.5\linewidth}p{0.45\linewidth}@{}}
\toprule
Markup & Meaning \\
\midrule
\cmd{usepackage\{zitex\}}                         & load the package (needs XeLaTeX) \\
\cmd{zi[style=S, lang=L, reading=R]\{字\}}         & general form; \texttt{S} is one of tree, matrix, formula, math \\
\cmd{zitree}, \cmd{zimatrix}, \cmd{ziformula}, \cmd{zimath} & shortcuts for the four styles \\
\cmd{zitexsetlang\{L\}}, \cmd{zitexsetreading\{R\}}  & document-wide defaults \\
\cmd{zitexsetpath\{dir\}}                          & where the snippets are (default \texttt{zitex-out}) \\
\bottomrule
\end{tabular}
\end{center}

\end{document}
